\chapter{Variational Principles and Hamiltonian Formulation}
\label{ch:variational}

The AKSZ construction of the previous chapter provides a BV
master action $S_{\mathrm{RSVP}}$ on the graded space of superfields.
In this chapter we extract the classical (degree~$0$) variational
principle, identify the induced Euler--Lagrange equations in ordinary
fields $(\Phi,\vec v,S)$, and construct the classical Hamiltonian
formulation on the derived mapping stack.
This leads to a well-defined notion of conserved currents and prepares
the ground for the analytic theory developed in Chapters~9--11.



\section{The Classical Variational Functional}
\label{sec:variational-functional}

Let $M$ be a smooth, oriented $4$--manifold without boundary (the
boundary case will be handled in the BV--BFV formalism).
Restrict the AKSZ action $S_{\mathrm{RSVP}}$ to the classical (ghost
number~$0$) sector. Denote by
\[
(\Phi,\vec v,S)\in \Gamma(M,\mathcal{E})
\]
the physical lamphrodynamic fields.
Then the classical action takes the form
\begin{equation}
\label{eq:classical-action}
\mathcal{S}[\Phi,\vec v,S]
=\int_{M}\mathcal{L}(\Phi,\vec v,S,\partial\Phi,\partial\vec v,
\partial S)\,d^{4}x,
\end{equation}
where the Lagrangian density $\mathcal{L}$ is obtained by truncating
the AKSZ superfield expansion and selecting the degree~$0$ component
(including nonlinear terms arising from the $(-1)$--shifted
symplectic structure).

\begin{remark}
The explicit coordinate expression of $\mathcal{L}$ is cumbersome due
to the nonlinear couplings between $\Phi,\vec v$ and $S$. However, the
structural properties---gauge invariances, conservation laws, and
entropy-flux terms---follow from the $(-1)$--shifted geometry rather
than from the coordinate expansion.
\end{remark}



\section{Euler--Lagrange Equations}
\label{sec:EL-equations}

The Euler--Lagrange equations are obtained by taking first variations
of \eqref{eq:classical-action}:
\begin{align}
\frac{\delta\mathcal{S}}{\delta\Phi}
&=
\partial_{\mu}\left(
\frac{\partial\mathcal{L}}{\partial(\partial_{\mu}\Phi)}
\right)
-
\frac{\partial\mathcal{L}}{\partial\Phi}
=0,
\label{eq:EL-Phi}
\\[0.5em]
\frac{\delta\mathcal{S}}{\delta v_{\nu}}
&=
\partial_{\mu}\left(
\frac{\partial\mathcal{L}}{\partial(\partial_{\mu}v_{\nu})}
\right)
-
\frac{\partial\mathcal{L}}{\partial v_{\nu}}
=0,
\label{eq:EL-v}
\\[0.5em]
\frac{\delta\mathcal{S}}{\delta S}
&=
\partial_{\mu}\left(
\frac{\partial\mathcal{L}}{\partial(\partial_{\mu}S)}
\right)
-
\frac{\partial\mathcal{L}}{\partial S}
=0.
\label{eq:EL-S}
\end{align}

\begin{theorem}[Equivalence with AKSZ Euler--Lagrange Equations]
\label{thm:EL-equivalence}
The equations \eqref{eq:EL-Phi}--\eqref{eq:EL-S} coincide with the
classical Euler--Lagrange equations obtained from the AKSZ action
$S_{\mathrm{RSVP}}$ by projecting to degree~$0$.
\end{theorem}

\begin{proof}
The AKSZ construction guarantees that the BV Euler--Lagrange
operator $\delta S_{\mathrm{RSVP}}/\delta\mathsf{Field}=0$ restricts
to the classical sector by ghost number.
Truncating to ghost number~$0$ and using the fact that the BV bracket
reduces to the classical Poisson bracket on physical fields yields
\eqref{eq:EL-Phi}--\eqref{eq:EL-S}.
\end{proof}



\section{Hamiltonian Formulation on the Derived Mapping Stack}
\label{sec:Hamiltonian}

Let $\mathrm{Map}(M,\mathrm{Plenum})$ be the classical mapping space
of $M$ into the derived target $\mathrm{Plenum}$,
and let $\omega_{(-1)}$ denote the $(-1)$--shifted symplectic form on
$\mathrm{Plenum}$.
In the AKSZ formalism, the Hamiltonian structure on the space of
superfields descends (after truncation) to a classical Hamiltonian
system on $\mathrm{Map}(M,\mathrm{Plenum})$ endowed with the induced
(shifted) 2--form.

\begin{proposition}[Classical Hamiltonian Structure]
\label{prop:Hamiltonian}
The classical space of fields $\mathrm{Map}(M,\mathrm{Plenum})$
carries a well-defined Hamiltonian structure $(\omega,H)$ where
\[
H=\int_{M} \mathcal{H}(\Phi,\vec v,S,
\partial\Phi,\partial\vec v,\partial S)\,d^{4}x,
\]
and the Hamiltonian density~$\mathcal{H}$ is obtained by classical
Legendre transform of the Lagrangian density~$\mathcal{L}$.
\end{proposition}

\begin{proof}[Proof (Sketch)]
The $(-1)$--shifted symplectic form on $\mathrm{Plenum}$ induces a
$0$--shifted (classical) symplectic form on the mapping space.
The Legendre transform is defined fiberwise on the classical field
bundle $\mathcal{E}\to M$.
Standard AKSZ arguments show that the resulting Hamiltonian generates
the same equations of motion as \eqref{eq:EL-Phi}--\eqref{eq:EL-S}.
\end{proof}

\begin{remark}[Odd vs Even Poisson Brackets]
The BV bracket on the AKSZ superfields is odd (degree $+1$),
whereas the induced bracket on the classical fields is even.
This provides a natural separation between quantum BV information
(ghost/antifield sector) and classical conservation laws.
\end{remark}

\section{Noether Currents and Conserved Quantities}
\label{sec:noether}

Let $\mathcal{G}$ be a Lie group of symmetries acting on the classical
fields $(\Phi,\vec v,S)$ via smooth bundle automorphisms
compatible with the lamphrodynamic constraint.
Write $\rho:\mathfrak{g}\to\Gamma(T\mathcal{E})$ for the induced Lie
algebra action on fields.

\begin{definition}[Infinitesimal Symmetry]
A vector field $X\in\Gamma(T\mathcal{E})$ is an infinitesimal
symmetry if the Lie derivative of the Lagrangian density satisfies
\[
\mathcal{L}_{X}\mathcal{L} = \partial_{\mu} B^{\mu}
\]
for some current $B^{\mu}$.
\end{definition}

Standard variational arguments yield:

\begin{proposition}[Noether Current]
\label{prop:noether-current}
Let $X$ be an infinitesimal symmetry associated with $\xi\in\mathfrak{g}$.
Then the Noether current is
\[
J^{\mu}_{\xi}
= \frac{\partial\mathcal{L}}{\partial(\partial_{\mu}\Phi)}X(\Phi)
 +\frac{\partial\mathcal{L}}{\partial(\partial_{\mu}v_{\nu})}X(v_{\nu})
 +\frac{\partial\mathcal{L}}{\partial(\partial_{\mu}S)}X(S)
 - B^{\mu},
\]
and is conserved along classical solutions:
$\partial_{\mu}J^{\mu}_{\xi}=0$.
\end{proposition}

\begin{remark}
In the BV formulation of Chapter~7, the same current arises as the
degree--$0$ component of the BV charge associated to the ghost--$\xi$.
Hence classical conservation laws are shadows of BV identities.
\end{remark}



\section{Hamiltonian Vector Fields and Poisson Structure}
\label{sec:ham-vfields}

Let $\omega$ denote the classical symplectic form on
$\mathrm{Map}(M,\mathrm{Plenum})$.
For any classical observable $\mathcal{O}$, define the Hamiltonian
vector field $X_{\mathcal{O}}$ by the relation
\[
\iota_{X_{\mathcal{O}}}\omega = \delta\mathcal{O}.
\]
Provided $\omega$ is weakly nondegenerate on the classical sector,
this determines $X_{\mathcal{O}}$ formally and induces a Poisson
bracket
\[
\{\mathcal{O}_{1},\mathcal{O}_{2}\}
:= X_{\mathcal{O}_{1}}(\mathcal{O}_{2}),
\]
which restricts to field polynomials under classical truncation.

\begin{remark}
Weak nondegeneracy follows from nondegeneracy of the $(-1)$--shifted
symplectic form on the target stack $\mathrm{Plenum}$ and the AKSZ
descent; a detailed discussion will appear in Appendix~E.
\end{remark}



\section{Energy, Momentum, and Entropic Contributions}
\label{sec:energy-momentum}

Energy--momentum for lamphrodynamic fields requires special care
because $S$ contributes an entropic stress--energy component not
present in standard relativistic field theories.

\begin{definition}[Lamphrodynamic Stress--Energy Tensor]
\label{def:stress}
Define the symmetric tensor
\[
T_{\mu\nu}
:=
\frac{\partial\mathcal{L}}{\partial(\partial^{\mu}\Phi)}
 \partial_{\nu}\Phi
+\frac{\partial\mathcal{L}}{\partial(\partial^{\mu}v_{\rho})}
 \partial_{\nu} v_{\rho}
+\frac{\partial\mathcal{L}}{\partial(\partial^{\mu}S)}
 \partial_{\nu}S
- g_{\mu\nu}\mathcal{L},
\]
where $g_{\mu\nu}$ is the background Lorentzian metric on $M$.
\end{definition}

\begin{proposition}[Energy Balance Identity]
\label{prop:energy-balance}
Along classical solutions of the Euler--Lagrange equations
\eqref{eq:EL-Phi}--\eqref{eq:EL-S},
the stress--energy tensor satisfies
\[
\nabla^{\mu}T_{\mu\nu}
= \mathscr{E}_{\nu}(\Phi,\vec v,S),
\]
where the entropic source term
$\mathscr{E}_{\nu}$ vanishes for isentropic flows and encodes
entropy--flux contributions in general.
\end{proposition}

\begin{remark}
The presence of $\mathscr{E}_{\nu}$ is physically crucial: it
generates effective geometric backreaction in emergent gravity
(Volume~II, Chapter~1).
\end{remark}



\section{Hamiltonian Energy Estimates and PDE Well--Posedness}
\label{sec:ham-energy}

The classical Hamiltonian
$H=\int_{M}\mathcal{H}\,d^{4}x$
provides the fundamental energy control for the analytic theory of
Chapters~9--11.
The following preparatory result connects the formal Hamiltonian
with Sobolev control.

\begin{theorem}[Hamiltonian Energy Estimate]
\label{thm:ham-estimate}
Let $(\Phi,\vec v,S)$ be a classical solution of the lamphrodynamic
field equations on a globally hyperbolic spacetime $(M,g)$.
Assume initial data in $H^{k}$, $k\ge2$,
and suppose the entropic term satisfies a coercivity condition
$\mathscr{E}_{0}\le C|\nabla S|^{2}$.
Then for sufficiently small time $T>0$,
\[
\sup_{t\in[0,T]}
\bigl( \|\Phi(t)\|_{H^{k}}
     +\|\vec v(t)\|_{H^{k}}
     +\|S(t)\|_{H^{k}}
\bigr)
\;\le\;
C_{0} \exp(CT)\,
\bigl( \|\Phi_{0}\|_{H^{k}}
     +\|\vec v_{0}\|_{H^{k}}
     +\|S_{0}\|_{H^{k}} \bigr),
\]
where $C_{0},C$ depend on the geometry of $M$.
\end{theorem}

\begin{proof}[Proof (Sketch)]
The Hamiltonian controls the $H^{k}$--norms modulo entropic terms.
Coercivity of $\mathscr{E}_{0}$ and Gr\"onwall's inequality yield the
exponential bound.
A detailed PDE proof is given in Chapter~10.
\end{proof}



\section{Summary and Forward Link}
\label{sec:summary}

In summary, the classical limit of the AKSZ--RSVP theory yields:
\begin{enumerate}
\item a well-defined variational principle,
\item Euler--Lagrange equations equivalent to the classical AKSZ equations,
\item a Hamiltonian formalism on the derived mapping stack,
\item Noether currents and entropic stress--energy,
\item energy estimates sufficient for local well-posedness.
\end{enumerate}
These results serve as the bridge between the geometric BV framework
(Chapters~4--7) and the analytic theory (Chapters~9--11),
where energy estimates and Sobolev bounds play a central role.

